\subsection{導入：なぜコンピュータ科学の基礎が必要か}

プログラマは、与えられたアルゴリズムを命令列へ変換し、コンパイル、リンク、ロード、実行を通じて出力を得ることに重点を置く。一方、ソフトウェア技術者は、要求、アーキテクチャ、設計、性能、テスト、受け入れ、保守、工程を含む広い判断を行う。したがって、ソフトウェア技術者は、単にプログラムを書くだけでなく、なぜその構成が妥当なのかを説明できなければならない。

本章でいうコンピューティング基礎は、ソフトウェア技術者が高品質なソフトウェアを設計、構築、配備、保守するために必要なコンピュータ科学の基本概念である。以降では、個々の技術を丸暗記するのではなく、設計判断の材料として理解することを重視する。

\begin{pointbox}{重要}
コードは最終成果物の一部であるが、設計判断はコードより前に始まり、運用・保守まで影響する。データ構造、計算量、基本ソフトウェア、ネットワーク、データベースの知識は、後工程の手戻りを減らすための共通語である。

\end{pointbox}
